\section{Test A: out-of-representation}
\label{sec:test_a_oor}

\subsection{Goal}
To avoid the objection that FBOHN reads its own label, a label independent of BOHN/FBOHN features was constructed:
\[
y=\mathbf{1}\!\left[
\sin(x_3x_8)+\cos(x_5x_{11})+0.75x_1x_7x_{14}+0.5x_2x_{13}
-0.35x_4x_9+0.25x_0x_6-0.20x_{10}x_{17}+\eta>\tau
\right].
\]

\subsection{Results}
\begin{center}
\begin{tabular}{rrrr}
\toprule
Noise & RAW & BOHN & FBOHN\\
\midrule
0.0 & 0.550 & 0.605 & 0.622\\
0.1 & 0.545 & 0.600 & 0.626\\
0.3 & 0.543 & 0.589 & 0.615\\
0.5 & 0.539 & 0.584 & 0.609\\
\bottomrule
\end{tabular}
\end{center}

\subsection{Analysis}
The FBOHN effect does not disappear after removing the direct dependence of the label on FBOHN features. The advantage is smaller than in v1 but stable: approximately $7$--$8$ percentage points over RAW and approximately $2$--$3$ points over BOHN.

\subsection{Conclusion}
\[
\boxed{\text{FBOHN provides a real geometric advantage also outside its own label class.}}
\]
This is a methodologically more important result than RMIG-FBOHN v1, because it demonstrates the utility of RMIG geometry as an inductive bias.
